% 本文件由 LaTeX 排版 Agent 根据有效 translation.json 自动生成。
% author=Augustine of Hippo; work_id=1304; title=论信德与信经

\chapter{论信德与信经}

% structure: p0001 source=h1 target=omitted reason=page-title-already-rendered-by-parent

在一次全体北非主教会议于\Place{希波勒吉乌斯}召开时发表的演讲。\par

% structure: p0003 source=h2 target=section reason=relative-heading-rank; base=h2

\section{引言说明}

本著作的成书背景与日期，在奥古斯丁\WorkTitle{纠正}第一卷第十七章所述的一段陈述中有所说明。\par

由此我们得知，它最初是奥古斯丁——当时仅是一位司铎——应聚集于希波勒吉乌斯的主教会议的请求，在公开场合发表的演讲；随后应朋友们的要求，以书籍形式出版。北非教会全体大会，即于此地——今\Place{阿尔及尔}境内的\Place{博纳}——召开，时值公元393年；该会议因通过了一项坚决的抗议而具有历史重要性，该抗议反对当时在教会中赋予宗主教以其他地位的立场，并反对授予教会最高官员任何比单纯\Quote{首座主教}（\OriginalTerm{首座主教}{primæ sedis episcopus}）更权威或更教师性质的称号。\par

该著作是对所谓宗徒信经各条款的阐释。关于天主圣三中三位相互关系的探讨最为详尽；与此相关，特别是在运用存在、知识和爱这些类比时，以及在对这些类比及取自人类经验的其他例子的某些应用所提出的告诫中，我们看到了一些也在\WorkTitle{天主之城}、\WorkTitle{论天主圣三}及其教义著作中重复出现的观点。\par

\WorkTitle{纠正}中提到的段落如下：大约在同一时期，在主教们——他们授权我这样做——面前，并在希波勒吉乌斯举行的全体非洲主教全体会议上，我以司铎身份发表了一篇关于\WorkTitle{论信德与信经}的讨论。应一些与我们关系尤为密切、亲爱之人的一再强烈请求，我将这次研讨整理成书，其中主题本身成为论述的对象，尽管并未采用交给\OriginalTerm{候洗者}{competentes}背诵的那种特定言辞顺序。在此书中，当讨论肉身复活的问题时，我说：\Quote{身体将要复活，按照基督信仰，这信仰是不会欺骗人的。如果有人对此感到难以置信，那是因为他仅仅看肉身现在的样子，而没有考虑它将来会有的性质。因为在那天使般变化的时刻，它不再是血肉，而只是身体；}等等，以及我在那里就属地身体变成属天身体这一主题所说的其他陈述，正如宗徒从相同角度说话时所说：\BibleQuote{血肉不能继承天主的国}\BibleRef{格前 15:50}。但是如果有人对这些声明作了这样的理解，以至于认为属地的身体（如我们现在所拥有的）在复活中变成了属天的身体，以至于肢体或肉身的实质将不再存在，那么无疑他需要被纠正，要提醒他想一想主的身体：主在复活后以同样的肢体显现，不仅可以用眼睛看到，也可以用手摸到，并且通过祂所说的话更确实地证明祂拥有肉身：\BibleQuote{你们摸摸我，看看，因为鬼神没有骨肉，如同你们看我有的}\BibleRef{路 24:39}。因此，可以肯定的是，宗徒并没有否认肉身的实质将在天主的国中存在，而是用\Quote{血肉}这一名称指代那些按血肉生活的人，或指代肉身的明确败坏，这败坏在那时必定不再存在。因为在他说了\BibleQuote{血肉不能继承天主的国}之后，他接着说的——即\Quote{朽坏也不能继承不朽坏}——被正确地理解为是对前面陈述的解释。关于这个难以说服不信者的主题，任何阅读我最近著作\WorkTitle{论天主之城}的人，都会发现我以自己所能达到的最大谨慎进行了论述。此处所述之著作以此开篇：\Quote{既然经上记载：}等等。\par

% structure: p0008 source=h2 target=section reason=relative-heading-rank; base=h2

\section{第一章 论本著作的起源与目的}

1. \Emph{鉴于}\BibleQuote{义人因信德而生活}这一立场——它建立在宗徒教导的最坚实基础上，且已被写定；又鉴于这信德要求我们同时尽心和舌的职责——因为宗徒说：\BibleQuote{心里相信，就可成义；口里承认，就可获得救恩}——所以，我们理当牢记成义与救恩。因为我们注定要在永久的正义中将来为王，当然不能从现今的邪恶世界获得救恩，除非同时，在为邻人的救恩劳作时，我们也用口承认我们内心所怀的信德。我们必须以虔诚和谨慎的警觉为目标，设法防止该信德在我们内受到任何伤害——无论来自哪一方——通过异端者的欺诈伎俩（或狡猾的欺诈）。\par

但我们却在信经中拥有公教信仰，这信仰为信徒所知晓并记诵，以尽可能简洁的表达形式得以保存，正如情况所允许的那样；其（编纂）目的是：那些在基督内重生的人中的初学者和婴孩，尚未通过勤勉和属灵的运用以及对圣经的领悟而得到坚固的，应获得一个用寥寥数语概括的纲要，涵盖那些必须相信的事项，这些事项随后随着他们在谦卑与爱德的稳固坚定基础上前进，并上升至神圣教义（的高度）时，要以许多话语向他们解释。因此，正是在信经中如此编排的这些寥寥数语之下，大多数异端者试图隐藏他们的毒素；神圣的慈悲已经透过属灵的人——他们不仅配得接受并相信如此阐释的公教信仰，而且藉由主所赐的光照彻底理解并领悟它——抵挡了，并且仍在抵挡这些异端者。因为经上记载：\Quote{你们若不相信，必不能明白。}然而，对信仰的探讨有助于保护信经；但目的不在于将此种探讨本身代替信经，给那些接受天主恩宠的人去记诵和重复，而是藉由公教权威和更坚固的防御，保护信经中所保留的事项免受异端者的阴险攻击。\par

% structure: p0011 source=h2 target=section reason=relative-heading-rank; base=h2

\section{二、论天主及其独有之永恒}

2. 因为某些人试图让人接受这样的观点：天主圣父并非全能；并非他们胆敢明确宣称这一点，但在他们的传统中，他们被证明怀有并相信这样的概念。因为当他们肯定存在一种本性，全能天主并未创造它，但同时祂又用这本性塑造了这个世界（他们承认这世界被安排得井井有条），他们由此否认天主是全能的，以至于不相信祂能创造世界而不必为了建造世界而使用另一种先前已存在、且非祂自己所造的本性。如此，他们根据对工匠、建筑者和各种艺人的肉体习惯，这些人若无已准备好的材料的帮助，就无法实现自己技艺的效果。因此，这些人也同样理解世界的创造者并非全能，如果祂不能在没有另一非祂所造的本性作为材料的帮助下塑造这个世界的话。或者，如果他们确实承认天主——世界的创造者——是全能，那么理所当然，他们必须也承认祂从无中创造了祂所造的一切。因为，承认祂是全能，就不能有任何事物不是祂的创造物。因为虽然祂从某物中造出某物，如用泥土造人，但祂当然没有用任何非祂所造之物来造任何对象；因为产生泥土的大地，祂是从无中造的。即使祂用某种材料造了天和地本身，即宇宙及其中的万物，正如经上所记：\Quote{你从看不见的材料，或从无定形的材料，创造了世界}，我们绝无必要相信，这宇宙由之而成的材料本身，尽管可能是无定形的，尽管可能是看不见的，无论其存在方式如何，都可能自己存在，仿佛它是与天主\OriginalTerm{同永恒}{co-eternal}且\OriginalTerm{同永年}{co-eval}的。但是，无论它拥有何种存在方式，即它以某种方式存在，并能够接受不同事物的\OriginalTerm{形式}{form}，这并非出自它自己，而是出自全能天主之手，因祂的仁慈，一切事物存在——不仅是已经成形的对象，而且是一切\OriginalTerm{可成形的}{formable}对象。此外，\OriginalTerm{成形者}{formed}与\OriginalTerm{可成形者}{formable}之间的区别是：成形者已经接受了形式，而可成形者能够接受同样的形式。但那位赋予对象形式的同一者，也赋予其可成形的能力。因为万有最美的不变形式出于祂并在祂内；因此祂自己是独一的，祂将可能性赋予一切事物，不仅使它们实际美丽，而且使它们能够美丽。因此，我们完全有理由相信天主从无中创造了万物。因为，即使世界是由某种材料造的，这材料本身也是从无中造的；因此，按照天主最有序的恩赐，首先进入的是接受形式的能力，然后一切已成形的事物才得以形成。我们说这些，是为了使人不要认为圣经的言语彼此矛盾，因为一方面记着天主从无中造了万物，另一方面说世界是从无定形的材料造的。\par

3. 因此，正如我们相信全能天主圣父，我们应当坚持这样的观点：没有任何受造物不是由全能者所创造的。既然祂藉着圣言创造了万物，这圣言也被称为真理、德能和天主的智慧——同样，在许多其他称号下，我们的主耶稣基督，被推荐给我们的信德，也呈现于我们的理性领会，即我们的救赎者与统治者，天主子；因为那藉以万有得以建立的圣言，不可能由任何其他者所生，而只能由那藉着祂作为工具建立万有的那位所生——\par

% structure: p0014 source=h2 target=section reason=relative-heading-rank; base=h2

\section{三、论天主子及其作为圣言的特殊称号}

— 鉴于情况如此，我重复说，我们也相信\Person{耶稣基督}，天主子，圣父的独生子，也就是说，他的独生子，我们的主。然而，我们不应仅仅按我们设想自己话语的意义来理解这圣言，那些话语是由声音和口发出，划破空气而消散，持续的时间不长于其声音的延续。因为那圣言恒久不变：关于这圣言本身，当论及智慧时曾说道：\BibleQuote{她安于自身，使万物更新。}再者，他被称为圣父的圣言，其原因在于圣父藉着他而被人认识。因此，正如我们说话的本意在于，当我们说真话时，藉着我们的话语使我们的心意被听者知晓，并且我们内心秘密携带的一切，都可通过这类标记表达出来，以便另一个人能明智地理解；同样，天主圣父所生的这智慧，最恰当地被称为他的圣言，因为最隐秘的圣父藉着他被有资格的心灵所认识。\par

4. 现在，我们的心智与我们试图表达该心智的话语之间有很大的区别。我们确实不是生出可理解的话语，而是形成它们；在形成它们时，身体是底层的材料。然而，心智与身体之间有最大的区别。但天主在生出圣言时，生出了他自己所是的。他既不是从无中生出，也不是从任何已经制成和奠定的材料中生出；而是从他自身生出了他自己所是的。因为我们在说话时也以此为目标（我们将会看到），如果我们仔细考虑我们意志的倾向；不是当我们说谎时，而是当我们说真话时。因为那时我们努力的目的不就是，如果可能的话，将我们自己的心智带到听者的心智中，以便他能领会并彻底洞察它；这样我们确实可以安于自身，不退出自身，同时发出一个标记，使得藉着它，关于我们的知识能在另一个人身上产生；这样，就赐予我们的能力而言，另一个心智如同被心智发出，藉此它能够显露自身？我们这样做，尝试通过话语、单纯的声音、面容、身体的姿态——如此多的方式，实在渴望使内在的东西显明；因为我们无法完全发出这类东西本身。因此，说话者的心智无法完全被认识；因此也为谎言打开了空间。另一方面，天主圣父，他既有意愿也有能力以最大的真理向预定认识他的心智宣告自己，为了如此宣告自己，他生出了这圣言，就是那生下者自己所是的；这位圣言也被称为他的能力和智慧，因为藉着他，他做了万有，并有序地安排了它们；为此，有这些话论及他：\BibleQuote{她（智慧）从一端猛烈地达到另一端，并温柔地安排一切。}\par

% structure: p0017 source=h2 target=section reason=relative-heading-rank; base=h2

\section{第四章　论天主子既非由圣父所造，亦非小于圣父，及其降生成人}

5. 因此，天主的独生子既非由圣父所造；因为，照一位福音作者的话说：\BibleQuote{万物都是藉着他造的}；也非瞬间受生；因为永远智慧的天主，与他自己的永恒智慧常在；也非与圣父不平等，也就是说，在任何方面小于他；因为一位宗徒也这样说：\BibleQuote{他虽具有天主的形体，却没有以自己与天主同等为应当把持不放的。}因此，按此大公信仰，一方面，那些主张圣子与圣父是同一位的，被排除；因为（显然）这圣言不可能\Emph{与}天主同在，若他不是与天主圣父同在，并且（同样明显）\Emph{独一者}不与任何人\Emph{同等}。另一方面，那些主张圣子是一个受造物，尽管不像其余的受造物那样，也同样被排除。因为无论他们宣称受造物有多大，若它是受造物，它就被塑造和制造。因为\Emph{塑造}和\Emph{创造}意思相同；虽然在拉丁语的用法中，\Emph{创造}一词有时被用来代替更准确的说法\Emph{生}。但希腊语作了区分。因为我们称那个为\OriginalTerm{受造物}{creatura}，而他们称之为\OriginalTerm{受造物}{\GreekText{κτίσμα}}或\OriginalTerm{受造物}{\GreekText{κτίσις}}；当我们想无歧义地说话时，我们不用\OriginalTerm{创造}{creare}一词，而用\OriginalTerm{形成/建立}{condere}一词。所以，如果圣子是受造物，无论他多么伟大，他都是被造的。但我们相信的是他，\OriginalTerm{万物}{omnia}都是藉着他造的，而不是他，\OriginalTerm{其余的事物}{cetera}是藉着他造的。因为在此我们不能再以其他任何意义理解\Emph{万物}一词，只能理解为所有被造之物。\par

6. 然而，正如\BibleQuote{圣言成了血肉，居住在我们中间}，那由天主所生的同一智慧，也屈尊就卑，在人类中被造成。正如这句话所指涉的：\BibleQuote{主在祂道路的开端创造了我}。因为祂道路的开端，就是教会的元首，即穿上了人性的基督（\OriginalTerm{穿上了人性}{homine indutus}），藉着祂，我们得以获得生活的楷模，即一条通往天主的稳妥道路。我们因骄傲而堕落，除了谦卑，别无回归之路，正如对我们的始祖所说的：\BibleQuote{尝了，你们就会如同天主一样}。因此，我们的恢复者亲自认为，以自身展示这一谦卑的榜样是适宜的，祂\BibleQuote{虽具有天主的形体，却没有把持自己与天主同等，反而空虚自己，取了奴仆的形体}，为要使祂，那万有藉以造成的圣言，在祂道路的开端成为人。因此，就祂是独生子而言，祂没有兄弟；但就祂是首生子而言，祂认为配得将兄弟的称号赐给那些，在祂的卓越之后并藉着祂的卓越，藉着嗣子的义子身份，在恩宠中重生的人，这是按宗徒教导所传授给我们的真理。因此，\OriginalTerm{按本性的圣子}{naturalis filius}是由圣父的本体所生，唯一如此生的，祂所是的就是父所是的，天主出自天主，光出自光。而我们，不是本性的光，而是被那光照亮，使我们能在智慧中发光。因为正如经上说：\BibleQuote{那才是真光，照亮一切来到世上的人}。因此，我们在永恒之事的信德上，加上了我们主在时间中的安排，祂认为配得为我们承担并施行以拯救我们。因为就祂是天主的独生子而言，不能说祂\Emph{曾是}或\Emph{将是}，而只能说祂\Emph{是}；因为一方面，那\Emph{曾是的}，\Emph{现在不再是}；另一方面，那\Emph{将是的}，\Emph{尚未是}。因此，祂是不变的，不受时间条件和变化的制约。我认为，正是因为这一考虑，祂将那种名号启示给了祂的仆人梅瑟。因为当梅瑟问祂，如果他所被派遣的人民藐视他，他应该说谁打发了他时，祂这样说：\BibleQuote{我是自有者}。此后，祂又加了这句话：\BibleQuote{你要对以色列子民这样说：\OriginalTerm{那自有的}{Qui est}打发我到你们这里来}。\par

7. 由此，我相信，属神的人现在可以清楚看到，不可能存在任何与天主相反的本性。因为如果祂\Emph{是}——这个词只能恰当地用于天主（因为那真正\Emph{是}的，保持不变；凡变化的，曾是它现在所不是的，并将是它尚未是的）——那么天主就没有任何与祂自己相反的东西。如果有人问我们：什么与白相反？我们会回答：黑；问什么与热相反？回答：冷；问什么与快相反？回答：慢；所有类似的问题我们都会这样回答。然而，当问及\Quote{那\Emph{是}的}什么相反时，正确的回答是：\Quote{那\Emph{不是的}}。\par

8. 然而，正如我所说，在时间性的安排中，为我们的救恩和恢复，并因天主的良善在其中运作，我们可变化的本性被那不变的天主智慧所摄取，因此我们加上了为我们的益处而实行的救恩性时间事件的信德，即相信\Emph{由圣神藉童贞玛利亚所生的天主子}。因为藉着天主的恩赐，即圣神，如此伟大的天主以其极大的谦卑赐予我们，祂认为配得在童贞女的子宫中摄取\OriginalTerm{人的全部本性}{totum hominem}，居住在有形身体中，使之毫发\OriginalTerm{无损}{integrum}，并在离开时也毫无损伤。这个时间性的安排多次被异端者狡猾地攻击。但如果任何人掌握了公教信德，相信天主圣言摄取了人的全部本性，即身体、灵魂和心神，他就有足够的防御来对付那些异端者。因为，既然那个摄取是为了我们的救恩，人就必须警惕，不要相信我们本性中有某些部分与那个摄取无关，否则那些部分也可能与救恩无关。而既然除了肢体形状——这是根据种类差异赋予不同生物的形式——人几乎不与牲畜相区别，除非是拥有\OriginalTerm{理性精神}{rationali spiritu}，也称为\OriginalTerm{理智}{mens}，那么那信德——即认为天主的智慧摄取了我们的那一部分（我们与牲畜共享的部分），却没有摄取那被智慧之光所照耀的、人所特有的部分——如何能是正确的呢？\par

9. 此外，那些否认我们的主耶稣基督在玛利亚那里有过地上的母亲的人，也当受憎恶；然而，那救恩的计画同时尊荣两性，即男性和女性，并清楚表明：不仅祂所取的那一性属于天主的眷顾，而且祂藉以取得另一性的那一性也是如此，因为祂承受了男人的本性（\OriginalTerm{承受男人的本性}{virum gerendo}），并且祂由女人所生。祂说这句话：\Quote{女人，这于我和你有什么关系？我的时刻尚未来到}\BibleRef{若2:4}，并无任何理由迫使我们否认主的母亲。相反，祂更提醒我们明白：就祂是天主而言，祂没有母亲，因为祂正准备在变水为酒的神迹中显示祂位格尊威的一部分（\OriginalTerm{其位格尊威}{cujus majestatis personam}）。但就祂被钉十字架而言，祂是因祂的人性而被钉；而那就是说话时尚未到来的时刻：\Quote{这于我和你有什么关系？我的时刻尚未来到}\BibleRef{若2:4}；也就是说，那时我才会承认你。因为在那时，当祂作为人被钉时，祂承认了祂\OriginalTerm{人性的母亲}{hominem matrem}，并\OriginalTerm{极其人道地}{humanissime}将她托付给最心爱的门徒照管。再者，我们也不应因有人向祂报告祂的母亲和弟兄们来了，祂回答说：\Quote{谁是我的母亲？谁是我的弟兄？}\BibleRef{玛12:48}等话而动摇。相反，这应教导我们：当父母妨碍我们向弟兄们传播天主圣言的服务时，我们不应承认他们。因为如果有人因祂说\Quote{谁是我的母亲？}而断定祂在地上没有母亲，那么他也必然被迫否认宗徒们有地上的父亲，因为祂曾这样吩咐他们：\Quote{不要称地上的人为你们的父亲，因为你们的父只有一位，就是天上的父。}\BibleRef{玛23:9}\par

10. 我们也不应让女人怀胎的念头损害这信德，以致似乎有必要拒绝我们主的这种降生，仅仅因为\OriginalTerm{卑微的人认为它卑贱}{sordidi sordidam putant}。因为宗徒的这些话最为真实：\Quote{天主的愚妄比人明智}\BibleRef{格前1:25}，\Quote{在纯洁的人看来，一切事都是纯洁的}\BibleRef{铎1:15}。因此，那些持这种看法的人应当思量：这太阳的光线——他们不仅不赞美它为天主的受造物，反而将其当作神来崇拜——遍及世界各地，穿过污水沟和各种可怕之物的恶臭，并按其本性在其中运作，却从未因所沾染的任何污秽而变得卑污；尽管可见的光在本质上与可见的污秽更为接近。那么，天主的圣言——祂既不是有形体的，也不是可见的——又怎么会从女性身体上沾染污秽呢？祂在女性身体里取得了带着灵魂和精神的人性，借着这肉身的进入，圣言的尊威以较不直接的方式与柔弱的人体结合。由此可见，天主的圣言绝不可能被人的身体玷污，因为连人的灵魂也不会被身体玷污。灵魂被身体玷污，不是当它统治身体并使之生活的时候，而只是当它贪恋身体内可朽坏之物的时候。但如果这些人想要避免灵魂的玷污，他们更应该逃避这些谎言和亵渎。\par

% structure: p0024 source=h2 target=section reason=relative-heading-rank; base=h2

\section{第五章——论基督的受难、埋葬与复活}

11. 我们的主为我们降生，这（相对而言）是小小的\OriginalTerm{卑屈}{humilitas}；此外，祂还认为适宜为必死的人而死。因为\Quote{祂贬抑自己，听命至死，且死在十字架上}\BibleRef{斐2:8}，以免我们中有人虽然不怕一般的死亡，却对某种人们认为最可耻的死法感到战栗。因此，我们信祂 \Emph{在般雀比拉多执政时被钉十字架、埋葬}。之所以要加上审判官的名字，是为了让人们知道时间。此外，当相信这埋葬时，也要想起那新坟墓，它预示祂将复活进入新生命，正如童贞女的子宫预示祂将降生一样。因为那坟墓里没有埋葬其他死人，无论是之前还是之后；同样，那子宫里也没有怀过任何可朽之物，无论是之前还是之后。\par

12. 我们也信祂 \Emph{第三日从死者中复活}，成为将要跟随祂的弟兄们的首生者，祂将他们召叫为天主的义子，并认为适宜使他们成为祂的同承继者和同享福分者。\par

% structure: p0027 source=h2 target=section reason=relative-heading-rank; base=h2

\section{第六章——论基督升天}

13. 我们也相信祂\Emph{升了天堂}，这福地同样也为我们所有人所应许，祂说：\BibleQuote{他们要像在天上的天使}，在那座为我们众人之母的城，那在诸天上的永恒\Place{耶路撒冷}。然而，某些异教徒或不虔敬的外邦人，或异端者，常常因我们相信一个属地的身体升入天堂而感到冒犯。外邦人却大都竭力用哲学家的论据说服我们，断言天堂里不可能有任何属地的存在。因为他们不认识我们的圣经，也不明白经上所说：\BibleQuote{播种的是属血肉的身体，复活的是属灵的身体。}因为这样说并不是指身体变成了灵而成为灵；正如目前，我们的身体被称为\OriginalTerm{属血肉的}{animale}，并没有变成灵魂而成为\OriginalTerm{灵魂}{anima}。所谓属灵的身体，是指身体如此服从于灵，以至于适于天上的居所，所有软弱和一切属地的瑕疵都被改变并转化为天上的纯洁和稳定。这改变正是宗徒同样所说的：\BibleQuote{我们都要复活，但我们并不都要改变。}而这一改变不是变得更坏，而是更好，同一位宗徒教导说：\BibleQuote{我们都要改变。}至于主的身体在天堂何处、如何存在，这是一个过于好奇和多余的问题，无须深究。我们只当相信它在天堂里。因为考察天堂的奥秘不属于我们这软弱之人的本分，但怀着崇高和尊贵的意念对待主身体的尊荣，却属于我们的信德。\par

% structure: p0029 source=h2 target=section reason=relative-heading-rank; base=h2

\section{第七章 论基督坐在圣父的右边}

14. 我们也相信祂\Emph{坐在圣父的右边}。但这并不是要使我们以为天主圣父仿佛为人的形状所限制，以致我们思想祂时，会联想到右边或左边。同样，当经文明说圣父坐时，我们也不可幻想祂是屈膝而坐，以免陷入那种不虔敬，即宗徒所诅咒的那些人\BibleQuote{将不可朽坏的天主的光荣变为可朽坏的人的形像}。因为基督徒在圣殿中为天主设立任何这样的像都是非法的；因此，在心中设立这样的像更为可恶，因为心才是真正的天主的圣殿，只要它洁除了地上的情欲和错误。因此，\Quote{在右边}这句话，我们必须理解为至福中的位置，那里有义德、平安和喜乐；正如山羊羔被放在左边，即苦难中，由于不义、劳苦和折磨。与此相应，当经上说天主\Quote{坐}时，这一表述并非指明肢体的姿态，而是审判权能，这尊威从未丧失，祂总是\OriginalTerm{按应得的赏罚}{digna dignis tribuendo}施予人应得的报应；尽管在最后的审判中，天主独生子、生活者和死者的审判者那毫无疑问的光辉，在人类中必将更显明地彰显。\par

% structure: p0031 source=h2 target=section reason=relative-heading-rank; base=h2

\section{第八章 论基督的审判来临}

15. 我们也相信，在最合宜的时候，祂\Emph{将从那里降来，审判生者死者}：无论这些词是指义人和罪人，还是指那些在祂降临时从地上找到的、在死亡之前仍然活着的人称为\OriginalTerm{生者}{quick}，而\OriginalTerm{死者}{dead}则指在祂来临时要复活的人。这一时期的安排不仅如祂作为天主所具有的那样，是现在，也是过去和将来。因为我们的主曾在地上，如今在天上，将来也要在祂的光荣中作为生者死者的审判者。因为祂还要再来，如祂升天一样，根据\BibleBook{宗徒}{大事录}的权威。因此，祂在\BibleBook{}{默示录}中按照这一时期的安排说话，其中记载：\BibleQuote{那今在、昔在、将来永远常存者这样说。}\par

% structure: p0033 source=h2 target=section reason=relative-heading-rank; base=h2

\section{第九章 论圣神与天主圣三的奥迹}

16. 因此，我们的主的神圣受生和祂的人性经世，既已被如此系统地安排并交付于信德，为了完善我们关于天主的信德，在我们的宣认中又加上了［关于］圣神［的道理］。圣神在本性上并不低于圣父和圣子，而是，可以说，\OriginalTerm{同体}{consubstantial}和\OriginalTerm{同永}{co-eternal}的；因为此圣三是一位天主，并非意味着圣父与圣子和圣神是同一个位格，而是意味着圣父是圣父，圣子是圣子，圣神是圣神；并且此圣三是一位天主，正如经上所载：\BibleQuote{以色列啊！请听：主你们的天主是唯一的天主。} 同时，若有人逐一询问我们，并问我们说：\Quote{圣父是天主吗？}我们应回答：\Quote{祂是天主。}若问圣子是否是天主，我们也要同样回答。若有人就圣神这样问我们，我们也不应回答说祂不是天主；但要谨防［然而］仅按应用于人的意义来接受这句话，就如经上所载：\Quote{你们是神。}因为凡由圣父藉圣子、通过圣神的恩赐所造所成的一切，按本性没有一个是神。因为正是这同一圣三，当宗徒说：\BibleQuote{万物都出于祂，依赖祂，归于祂。}时所指明的。因此，虽然当我们逐一被问及每一位［这些位格］时，我们回答被询问的那一位——无论是圣父、圣子或圣神——是天主，但没有人应该因此认为我们敬拜三位神。\par

17. 这些事涉及不可言喻的本性而被说出，并不奇怪，因为即使在我们用肉眼观察、凭肉体感官判断的对象中，也有类似的情况。试以关于泉源的询问为例，思考我们如何不能断言那泉源本身就是河流；当我们被问及河流时，我们也同样不能称它为泉源；同样，我们也无法将从泉源或河流中取出的水称为河流或泉源。然而，在这个三分体中，我们使用\Emph{水}这个名称［作为整体］；当分别询问其中每一个时，我们在每个情况下都回答该物是\Emph{水}。因为若我询问泉源中是否有水，回答是肯定的；若我们询问河流中是否有水，回答并无不同；关于上述取出的水，也不可能作其他回答：然而尽管如此，我们不把这些称为三种水，而是一种水。当然，同时必须谨防，没有人仅仅像他想像这可见的、有形的泉源、河流或水那样来想像那威能不可言喻的本体。因为就后者而言，目前存在于泉源中的水流入河流，并不停留在自身中；当它从河流或泉源进入所取的水中时，它并不永久停留在被取出的地方。因此，在这里，同一水可能有时被称为泉源，有时被称为河流，有时被称为所取的水。但在那个圣三中，我们已断言圣父不可能有时是圣子，有时是圣神：正如在一棵树中，根只是根，树干（\Emph{robur}）只是树干，我们不能称树枝为别的什么；因为称为根的东西不能称为树干和树枝；属于根的木头不能通过任何转移时而存在于根中，时而存在于树干中，时而又存在于树枝中，而只存在于根中；这一命名的规则稳固不变，因此根是木头，树干是木头，树枝是木头，然而这样说并非指三种木头，而只指一种。或者，如果这些物体有某种不同，以至于因它们在强度上的差异，可以毫无荒谬地被说成三种木头；至少所有人都承认前一例子的力证——即，如果用同一泉源装满三个杯子，它们当然可以被称为三个杯子，但不能被说成三种水，而只是作为一个整体的一种水。然而同时，当逐一询问关于各个杯子时，我们可以回答每个杯子中都有水；尽管在这种情况下，并没有发生我们所说的从泉源流入河流那样的转移。但这些\OriginalTerm{有形的例证}{corporalia exempla}被引用，并非因为它们与神圣本性相似，而是因为即使在有形之物中也存在的统一性，以便使人理解，三类对象不仅各自、而且一起获得一个单一的名称是完全可能的；这样，就没有人感到奇怪，并认为我们称圣父为天主、圣子为天主、圣神为天主，却又说在圣三中并非三位神，而是一位天主、一个本体，是荒谬的。\par

18. 确实，关于圣父和圣子这个主题，博学而有灵性的人们已在许多著作中进行了讨论，在这些讨论中，他们尽可能以人与人之间的方式，努力引入一种可以理解的解释，说明圣父如何位格上不与圣子同一，然而二者在本体上是一；说明圣父\OriginalTerm{个别地}{proprie}是什么，圣子是什么：即前者是生者，后者是受生者；前者非出自圣子，后者出自圣父；前者是后者的元始（Beginning），因此他也被称为基督的元首，尽管基督同样是元始，但不是出于圣父；此外，后者是前者的肖像，虽毫无相似之处，且\OriginalTerm{绝对地且无差别地平等}{omnino et indifferenter æqualis}。这些问题被那些在比我们目前更广阔的范围内寻求完整表述基督徒信仰的人更广泛地处理。因此，就祂是圣子而言，祂从圣父领受了\Emph{祂是}，而另一位（圣父）却没有从圣子领受这点；并且，就祂以不可言喻的仁慈、在时间性的安排中取了自己人性（\OriginalTerm{人}{hominem}）——即那可变的受造物，并借此被改变成更美好的——而言，在圣经中发现许多关于祂的陈述，这些陈述的表达方式给异端者的不虔诚理智提供了错误的契机，他们以教导的愿望先于理解的愿望，以致他们以为祂既不与圣父平等，也不与圣父同一实体。这样的陈述例如：\BibleQuote{因为父比我大}\BibleRef{若 14:28}、\BibleQuote{女人的头是男人，男人的头是基督，基督的头是天主}\BibleRef{格前 11:3}、\BibleQuote{那时，他自己也要屈服于那使万物屈服于他的那一位}\BibleRef{格前 15:28}、\BibleQuote{我往我父、也往你们的父、我的天主和你们的天主那里去}\BibleRef{若 20:17}，以及其他一些类似性质的陈述。所有这些陈述被给予位置，当然不是意在表示本性或实体的不平等；因为若如此理解，就会歪曲另一类陈述，如：\BibleQuote{我与父是一体}（\OriginalTerm{一体}{unum}）\BibleRef{若 10:30}、\BibleQuote{看见我的，就是看见了父}\BibleRef{若 14:9}、\BibleQuote{圣言就是天主}\BibleRef{若 1:1}，因为祂不是受造的，因为\BibleQuote{万物是藉着祂造成的}\BibleRef{若 1:3}、\BibleQuote{祂不以自己与天主同等为僭夺}\BibleRef{斐 2:6}，以及其他一切类似经文。但这些陈述被给予位置，部分是为了祂\OriginalTerm{所取人性的职务}{administrationem suscepti hominis}，根据这一职务，经上说：\BibleQuote{祂空虚了自己}\BibleRef{斐 2:7}，这并不是说那智慧改变了，因为它绝对不变；而是祂愿意以如此谦卑的方式向人显示自己。因此，我重复说，部分是为了这一职务，那些异端者赖以作出虚假指控的事情才如此记载；部分是为了考虑圣子\Emph{所是}的一切都归于圣父——因而也特别归给圣父这一点，即祂\OriginalTerm{与同一圣父平等或相等}{eidem Patri æqualis aut par est}，而圣父对任何人都不欠祂所是的一切。\par

19. 然而，关于圣神，博学而卓越的圣经研究者们迄今尚未进行足够充分或细致的讨论，以使我们能够获得对祂的\OriginalTerm{特殊个体性}{proprium}的明智认识；正是凭借这一特殊个体性，我们不能称祂为圣子或圣父，而只能称祂为圣神；只不过他们断言祂是天主的恩赐，好让我们相信天主不会赐予比祂自身低劣的恩赐。同时，他们坚持这一立场，即既不承认圣神像圣子那样由圣父所生（因为基督是唯一如此生的），也不承认祂由圣子所生，如同至高圣父的孙子；他们不仅不主张祂的所是无须归因于任何存在，反而承认祂的所是要归因于圣父，万有皆从祂而来；免得我们确立\OriginalTerm{两个无始的元始}{ne duo constituamus principia isne principio}，这将是极其错误和荒谬的论断，不属于公教信仰，而属于某些异端者的谬误。然而，有些人竟至于相信圣父与圣子的共融，也就是（可以这么说）他们的\OriginalTerm{神性}{deitatem}，希腊人称之为\OriginalTerm{神性}{\GreekText{θεότης}}，就是圣神；因此，既然圣父是天主，圣子是天主，那么这神性本身，即他们彼此联合于其中的神性——即圣父藉生圣子，圣子藉依附于圣父——就应当与生祂的那一位平等。这神性，他们同样理解为存在于这两位之间的彼此相爱与爱德，他们断言这神性被称为圣神。他们用许多圣经的证明来支持这一见解；其中我们可以举出诸如这样的经文：\BibleQuote{天主的爱藉着赐予我们的圣神，已倾注在我们心中}\BibleRef{罗 5:5}或其他许多类似含义的经文；他们也把立场建立在这样明确的事实上：我们是通过圣神与天主和好；因此，当祂被称为天主的恩赐时，他们便认为这足以表明天主的爱与圣神是同一的。因为我们唯有通过那使我们得以称为子女的爱，才能与祂和好：我们不再\BibleQuote{像仆人一样，处于恐惧之下}\BibleRef{罗 8:15}，因为\BibleQuote{爱一旦达到完美，就驱逐恐惧}\BibleRef{若一 4:18}；并且\BibleQuote{我们领受了使人自由的神，在这神内我们呼喊：阿爸，父啊！}\BibleRef{罗 8:15}。既然藉着爱得以和好并被召回友谊中，我们将能认识天主的一切奥秘，因此论到圣神说：\BibleQuote{祂要引导你们进入一切真理}\BibleRef{若 16:13}。同样的理由，宗徒们在圣神降临时所充满的宣讲真理的信心，也应正当地归因于爱；因为胆怯则归因于恐惧，而爱的完美排除恐惧。同样，圣神被称为天主的恩赐，因为没有人能享受他所认识的，除非他也爱它。然而，享受天主的智慧无非就是藉着爱依附于祂（\OriginalTerm{藉着爱依附}{ei dilectione cohærere}）。也没有人凭爱之外的方式持守他所领悟的；因此圣神被称为\OriginalTerm{圣德之神}{Spiritus Sanctus}，因为一切\OriginalTerm{被圣化}{sanciuntur}的事物都是为了永久而圣化，而\OriginalTerm{圣德}{sanctitatem}一词无疑源自\OriginalTerm{圣化}{a sanciendo}。然而，这种见解的支持者特别使用这一见证，即经文所载：\BibleQuote{从肉身生的就是肉身，从神生的就是神}\BibleRef{若 3:6}；\BibleQuote{因为天主是神}\BibleRef{若 4:24}。因为这里祂论及我们的重生，这重生按照亚当不是出于肉身，而是按照基督出于圣神。因此，如果在这段经文中，当它说\BibleQuote{因为天主是神}时提到了圣神，他们主张我们必须注意，它说的不是\Quote{因为神是天主}，而是\BibleQuote{因为天主是神}；如此，圣父和圣子的神性本身在这里被称为天主，那就是圣神。再加上宗徒若望提供的另一个见证，他说：\BibleQuote{因为天主是爱}\BibleRef{若一 4:8}。因为这里，同样地，他说的不是\Quote{爱是天主}，而是\BibleQuote{天主是爱}；如此，神性本身就被认为是爱。至于这种情形：在列举相互关联的对象时，如经上所说：\BibleQuote{一切都是你们的，你们是基督的，基督是天主的}\BibleRef{格前 3:22-23}，以及\BibleQuote{女人的头是男人，男人的头是基督，基督的头是天主}\BibleRef{格前 11:3}，其中并未提及圣神；他们肯定这不过是应用了一个原则：一般而言，联系本身通常不被列在相互关联的事物之中。因此，那些更仔细阅读的人也似乎在另一处经文中同样认出明确的天主圣三，即经上说：\BibleQuote{因为万物都出于祂，藉着祂，并在祂内}\BibleRef{罗 11:36}。\Quote{出于祂}，仿佛意指出于那一位，祂的存在不归因于任何人；\Quote{藉着祂}，仿佛意思是藉着中保；\Quote{在祂内}，仿佛是在那一位内，祂维系一切，即藉着联系而结合。\par

20. 那些认为我们所称的神性、爱或爱德这一共融并非实体的派别，反对这一意见。此外，他们要求按实体来阐述圣神；他们也不认为，除非爱是实体，\BibleQuote{天主是爱} 这句话就不可能被使用。确实，他们在这一点上受到有形事物习惯的影响。因为如果两个物体彼此连接，并置在一起，连接本身并不是一个物体：因为当这些原本连接的物体分离时，当然不再发现这样的连接；同时，它也不被理解为像那些物体本身那样离开并迁移。但这样的人应尽其所能洁净自己的心，以便能够看见：在天主的实体中，没有任何东西会意味着实体是一回事，而 \Emph{\OriginalTerm{依附于实体的东西}{aliud quod accidat substantiae}} 是另一回事，且不是实体；而凡能被视作在其中者，都是实体。然而，这些事容易说出并相信，但要看透它们本身如何，则除非藉着洁净的心，否则绝不可能。因此，无论这意见是否正确，或另有实情，信仰都应当毫不动摇地持守：我们应称圣父为天主，圣子为天主，圣神为天主，但不说三个天主，而视上述圣三为独一天主；同样，不说这些位格在本性上不同，而视他们为同一实体；并且进一步坚持：并非圣父有时是圣子，有时是圣神，而是圣父永远是圣父，圣子永远是圣子，圣神永远是圣神。关于不可见之事，我们也不应像有知识一样轻率地断言，而应谦虚地相信。因为除非藉着洁净的心，这些事是无法看见的；即使是今生 \BibleQuote{部分地}，正如经上所说，和 \BibleQuote{如同在镜子里} 看见它们的人，也不能保证听他说话的人也看见它们，如果他受心灵的不洁所阻碍。然而，\BibleQuote{心里洁净的人是有福的，因为他们要看见天主。} 这就是关于我们的创造者和更新者天主的信仰。\par

21. 然而，既然爱不仅被命令指向天主，当时说：\BibleQuote{你应全心、全灵、全意爱上主你的天主；} 也指向我们的邻人，因为他说：\BibleQuote{你应爱你的邻人如你自己；} 而且，此外，如果这信仰不包含一个可通过兄弟爱德的运作的聚会和人群，它就不那么丰硕——\par

% structure: p0040 source=h2 target=section reason=relative-heading-rank; base=h2

\section{第十章  论天主教会、罪过的赦免与肉身的复活}

——既然如此，我们相信 \Emph{圣教会}，即确凿无疑的 \Emph{天主教会}。因为异端者和分裂者都称他们的聚会为教会。但异端者由于对天主持有错误见解，损害了信仰本身；而分裂者则因邪恶的分裂脱离了兄弟友爱，尽管他们可能相信我们所信的一切。因此，异端者不属于爱天主的天主教会，分裂者也不属于同一教会，因为它爱邻人，从而乐意赦免邻人的罪过，因为它祈求那使我们与祂和好的天主赦免自己，祂除去一切旧物，召叫我们活出新生命。直到我们达至这新生命的圆满，我们不可能没有罪过。然而，这些罪过属于何种性质是重要的。\par

22. 我们不仅应当讨论罪过的区别，更应彻底相信，如果我们对他人的罪过严苛不肯原谅，那么我们自己的罪过就绝不会得到赦免。因此，我们也相信 \Emph{罪过的赦免}。\par

23. 既然人由三部分构成——即精神、灵魂和身体——而又常被说成两部分，因为灵魂常与精神一同被提及；因为同一位阶的理性部分，为野兽所无，被称为精神：我们内主要的部分是精神；其次，使我们与身体结合的生命称为灵魂；最后，身体本身，作为可见的，是我们内最后的部分。然而，这\BibleQuote{整个受造物}（\OriginalTerm{受造物}{creatura}），\BibleQuote{直到现在都一同叹息、一同受产痛}。尽管如此，祂已赐给它圣神的初果，因为受造物已相信了天主，且如今怀着善意。这精神也称为心智，关于此，一位宗徒这样说：\BibleQuote{我用心智服事天主的法律。}同一宗徒在另一处也这样表达：\BibleQuote{天主为我作证，我在我的精神中事奉祂。}此外，灵魂，当它还贪恋属肉体的美物时，被称为肉体。因为它的某一部分抵抗圣神，不是出于本性，而是出于罪恶的习惯；因此经上说：\BibleQuote{我用心智服事天主的法律，却用肉体服事罪的法律。}这习惯已因原祖的罪，按照有死的世代，变成了本性。因此也这样记载着：\BibleQuote{我们从前也生来就是义怒之子，}即报复之子，因此我们服事了罪的法律。然而，灵魂的本性，当它服从于自己的精神，且当它跟随那精神如同精神跟随天主时，就是完美的。因此\BibleQuote{属血气的人不接受天主圣神的事。}但是，灵魂不至于像精神服从天主得到真信德和善意那样迅速地服从精神而行动；有时它向下趋向属肉体、暂时之事的冲动，被更迟缓地约束。然而，既然这同一灵魂也被净化，并在精神（它是灵魂的头）的统治下获得自身本性的稳定——而这灵魂的头又有自己的头，就是基督——我们也不应绝望于身体也会恢复其自己适当的本性。但这当然不会像灵魂那样迅速实现；正如灵魂也不像精神那样迅速恢复。然而，在适当的季节，在最后的号角声中，它将要实现，那时\BibleQuote{死人要复活，成为不朽的，我们也要被改变。}因此我们也信\Emph{肉身的复活}，就是说，不仅那因肉情而被叫做肉体的灵魂得以恢复，而且这可见的、按本性称为肉体的身体也将如此恢复；灵魂得到\BibleQuote{肉体}这名称，不是出于本性，而是因为肉情：所以我说，这可见的、严格意义上的肉体，无疑应被相信注定要复活。因为宗徒保禄似乎用手指着它说：\BibleQuote{这可朽坏的必须穿上不可朽坏的。}因为当他说\Emph{这}时，他仿佛用手指指向它。现在，只有可见的东西才可能被手指指出来；因为灵魂也可能被称为可朽坏的，因为它被恶习所败坏。而当我们读到\BibleQuote{这必死的必须穿上不死的}时，同样指的是可见的肉体，因为经文仿佛不断地用手指指向它。同样，灵魂也可以被称为必死的，正如它因恶习而被指为可朽坏的。因为确实，\Quote{灵魂的死就是背弃天主；}这是它在乐园中的第一个罪，正如圣经所载。\par

24. 因此，根据基督信仰——这信仰绝不会欺骗人——身体必将复活。如果有人觉得这事难以置信，那是因为他只看到血肉当下的样子，而没有考虑它将来会是什么性质。因为在那个天使的变化时刻，它不再是血肉，而仅仅是身体。因为当宗徒谈到血肉时，他说：\BibleQuote{有牲畜的肉、鸟的肉、鱼的肉、爬虫的肉；也有属天的身体和属地的身体。}然而他在这里所说的不是\Quote{属天的血肉}，而是\Quote{属天的身体和属地的身体}。因为凡血肉也都是身体，但并非每个身体也都是血肉。首先，例如在属地体的情况下，木头是身体，但不是血肉。在人或是牲畜的情况下，我们既有身体也有血肉。另一方面，在属天体的情况下，没有血肉，只有那些单纯明亮的身体，宗徒称之为属神的，而有些人称之为以太的。因此，当他说：\BibleQuote{血肉不能继承天主的国}时，这与血肉的复活并不矛盾；而是这话预言了那将来之物的性质——那目前是血肉的，将来会成为怎样。如果有人拒绝相信血肉能够变成所指示的那类性质，就必须逐步引导他走向这个信仰。因为如果你问他，土是否能变成水，因为物质间的相近，他不会觉得难以置信。同样，如果你问，水是否能变成气，他也会回答这并不荒谬，因为元素彼此相近。而如果问及气，它是否能变成以太体，即属天体，这种相近性本身就使他立即相信其可能性。但是，如果他已经承认通过这些层次，土变成以太体是可能的，那么当加上天主的旨意——藉此人的身体能够在水上行走——他为什么拒绝相信同样的变化可以极快地发生，正是按所说的\BibleQuote{在一眨眼之间}，而不需这些层次，正如通常烟雾奇迹般地迅速变成火焰？因为我们的血肉确实来自土。但哲学家们——主要依据他们的论证才有人反对血肉的复活——尽管他们断言没有属地体能在天上存在，却承认任何一种身体都可以转化和改变成任何其他类型的身体。当这身体的复活完成后，从时间状态中解放出来，我们将要在无法言说的爱与坚定不移中，全然地享受\Emph{永生}，没有败坏。因为\BibleQuote{那时就要应验经上记载的话：死亡在胜利中被吞灭了。死亡啊！你的刺在哪里？死亡啊！你的争战在哪里？}\par

25. 这就是以简短话语在\WorkTitle{信经}中赐予基督徒初学者、让他们持守的信仰。这些简短话语为信友所熟知，目的在于藉着相信，他们得以服从天主；藉着服从，他们得以正当地生活；藉着正当地生活，他们得以洁净内心；藉着洁净的内心，他们得以明白所相信的。\par

\SourceRecord{Translated by S.D.F. Salmond.；Nicene and Post-Nicene Fathers, First Series；Vol. 3.；Edited by Philip Schaff.；Buffalo, NY: Christian Literature Publishing Co.,；1887.；<http://www.newadvent.org/fathers/1304.htm>.}
